\chapter{Consumer Welfare}\label{ch:b1:consumer-welfare}

From 1990 to 2024 the American consumer price index rose by a factor of
2.40, an average of $2.61\%$ a year. Over the same years the deflator of
household consumption, another price index for the same households, rose
by a factor of 2.07, $2.16\%$ a year. % figdata: us_growth
Social Security benefits are tied to the CPI-W, a sibling of this index,
and many pensions and wage contracts to fixed-basket indices like
it; % ledger: b1:cpi-uses
since 2017 US federal tax brackets have followed a chained index
instead. % ledger: b1:ccpiu-tax
If a fixed-basket index
overstates the rise in the cost of living by half a point a year, a
pension indexed to it gains about $18\%$ in real terms over 34 years. This
chapter defines what ``the cost of living'' means in the model of
\cref{ch:b1:consumer-choice}, measures the welfare effect of a price change
in money, and shows which way a fixed-basket index errs.

\begin{openingchart}{Two measures of US consumer prices, 1990--2024, 1990 = 100:
the consumer price index for all urban consumers (CPI-U, annual average), a
fixed-basket index whose weights are updated at intervals, and the
deflator of household final consumption from the national accounts, a
chained index. Scope and weights differ as well as the
formula.}{U.S. Bureau of Labor Statistics, consumer price index, all
items, U.S. city average; World Bank, World Development Indicators,
household final consumption expenditure in current and constant prices}
\begin{tikzpicture}
\begin{axis}[omtime, xlabel={Year}, ylabel={Index, 1990 = 100},
  xmin=1990, xmax=2024, ymin=90, ymax=250, width=0.72\linewidth,
  height=0.46\linewidth, legend pos=north west]
  \addplot+[omDef, very thick, no marks] table[x=year, y=cpi] {figdata/bachelor-1/07-consumer-welfare/out/us_indices.csv};
  \addplot+[omThm, thick, dashed, no marks] table[x=year, y=deflator] {figdata/bachelor-1/07-consumer-welfare/out/us_indices.csv};
  \legend{CPI-U, Consumption deflator}
\end{axis}
\end{tikzpicture}
\end{openingchart}

The two lines start together and end 33 index points apart: 240.01
against 206.89 in 2024. % figdata: us_growth
Part of the gap is about formula: a fixed basket ignores that households
buy less of what has become dear. Part is about scope: the deflator also covers
some spending made on households' behalf, such as health care paid by
employers and governments, which the CPI leaves out. The chapter's
theory isolates the first part.

\begin{recall}
A Laspeyres index prices a fixed base-period basket at current prices; a
Paasche index prices the current basket. A fixed basket overstates the rise
in the cost of living because it ignores substitution (the substitution
bias); a chained index updates the basket every period. \omterm{def:b1:consumer-welfare:consumer-surplus}{Consumer surplus} is
the area between the demand curve and the price. The \omterm{def:b1:income-and-substitution-effects:expenditure}{expenditure function}
$e(p, u)$ and Shephard's lemma, $\partial e/\partial p_i = h_i$, are in
\cref{ch:b1:income-and-substitution-effects}.
\end{recall}

\section{Money measures of a price change}\label{sec:b1:consumer-welfare:cv-ev}

A price rises from $p^0$ to $p^1$ (one price, the others fixed), and the
consumer's utility falls from $u^0 = v(p^0, m)$ to $u^1 = v(p^1, m)$. Utility
has no unit (\cref{prop:b1:preferences-and-utility:ordinal}), but income
does, and the \omterm{def:b1:income-and-substitution-effects:expenditure}{expenditure function} converts one into the other.

\begin{definition}[Compensating variation]\label{def:b1:consumer-welfare:cv}
\index{compensating variation}
The \emph{compensating variation} of the price change is
$\mathrm{CV} = e(p^1, u^0) - e(p^0, u^0) = e(p^1, u^0) - m$: the extra income that,
at the new prices, would restore the old utility.
\end{definition}

\begin{definition}[Equivalent variation]\label{def:b1:consumer-welfare:ev}
\index{equivalent variation}
The \emph{equivalent variation} is
$\mathrm{EV} = e(p^1, u^1) - e(p^0, u^1) = m - e(p^0, u^1)$: the loss of income
that, at the old prices, would hurt as much as the price change.
\end{definition}

For a price rise both are positive; for a price fall both are negative
(the consumer would pay $|\mathrm{CV}|$ to have the fall, or would need
$|\mathrm{EV}|$ more to be as well off without it).

\begin{proposition}[CV and EV as areas]\label{prop:b1:consumer-welfare:areas}
\index{compensating variation!as an area}
When only $p_i$ changes,
\[
\mathrm{CV} = \int_{p_i^0}^{p_i^1} h_i(p_i, u^0)\,\dd p_i ,\qquad
\mathrm{EV} = \int_{p_i^0}^{p_i^1} h_i(p_i, u^1)\,\dd p_i ,
\]
the areas to the left of the \omterm{def:b1:income-and-substitution-effects:compensated}{compensated demand} curves at the old and the
new utility, between the two prices. For a \omterm{def:b1:consumer-choice:normal-inferior}{normal good} and a price rise,
$\mathrm{CV} \ge \mathrm{EV}$.
\end{proposition}

\begin{proof}
For a fixed $u$, $e(p^1, u) - e(p^0, u) = \int \partial e/\partial p_i\,\dd p_i$
over $[p_i^0, p_i^1]$, and $\partial e/\partial p_i = h_i$ (Shephard's lemma,
\cref{prop:b1:income-and-substitution-effects:shephard}). With $u = u^0$ this is
CV, with $u = u^1$ EV. For a \omterm{def:b1:consumer-choice:normal-inferior}{normal good}, more utility means more of the good
at every price: $h_i(p_i, u) = x_i\bigl(p, e(p, u)\bigr)$ rises with $u$, since $e$
rises with $u$. A price rise gives $u^1 < u^0$, so $h_i(\cdot, u^0) \ge h_i(\cdot, u^1)$
and the first integral is the larger.
\end{proof}

\begin{omfigure}
\begin{tikzpicture}
\begin{axis}[omrecord, grid=none, xlabel={Quantity of good 1}, ylabel={Price $p_1$},
  xmin=0, xmax=720, ymin=0.5, ymax=2.5, width=0.6\linewidth, height=0.45\linewidth, samples=60,
  legend pos=north east]
  \addplot[fill=omDef!14, draw=none, domain=1:2, variable=\t, forget plot] ({353.55/sqrt(t)}, {t}) -- (axis cs:0,2) -- (axis cs:0,1) -- cycle;
  \addplot[shifted, domain=0.5:2.5, variable=\t] ({500/sqrt(t)}, {t});
  \addplot[demand, domain=0.5:2.5, variable=\t] ({500/t}, {t});
  \addplot[supply, domain=0.5:2.5, variable=\t] ({353.55/sqrt(t)}, {t});
  \addplot[guide, forget plot] coordinates {(0,1) (720,1)};
  \addplot[guide, forget plot] coordinates {(0,2) (720,2)};
  \legend{{$h_1(\cdot, u^0)$}, $x_1$, {$h_1(\cdot, u^1)$}}
  \node[font=\footnotesize] at (axis cs:110,1.5) {EV};
\end{axis}
\end{tikzpicture}

{\small $u = \sqrt{x_1x_2}$, income 1000, $p_2 = 1$, $p_1$ from 1 to 2. Between the
dotted price lines, the area left of the \omterm{def:b1:income-and-substitution-effects:compensated}{compensated demand} at the new
utility (red, shaded) is $\mathrm{EV} = 292.89$; left of the \omterm{def:b1:consumer-choice:demand-function}{Marshallian demand}
(blue) the consumer-surplus loss, 346.57; left of the \omterm{def:b1:income-and-substitution-effects:compensated}{compensated demand} at
the old utility (dashed) $\mathrm{CV} = 414.21$.}
\end{omfigure}

For the Cobb--Douglas consumer of the figure, $e(p, u) = 2u\sqrt{p_1p_2}$ and
$e$ is proportional to $\sqrt{p_1}$, so doubling $p_1$ raises the cost of any
utility by the factor $\sqrt 2$:
$\mathrm{CV} = 1000(\sqrt 2 - 1) = 414.21$ and
$\mathrm{EV} = 1000(1 - 1/\sqrt 2) = 292.89$.

\section{Consumer surplus}\label{sec:b1:consumer-welfare:surplus}

CV and EV need the \omterm{def:b1:income-and-substitution-effects:compensated}{compensated demand}, which is not observed. The demand
curve that is estimated from market data is the Marshallian one.

\begin{definition}[Consumer surplus]\label{def:b1:consumer-welfare:consumer-surplus}
\index{consumer surplus}
The \emph{consumer surplus} at price $p_i$ is the area under the
(Marshallian) demand curve above the price,
$\mathrm{CS}(p_i) = \int_{p_i}^{\infty} x_i(t, m)\,\dd t$ (with the integral
stopping at the price where demand reaches zero, if there is one). The
change in consumer surplus when the price rises from $p_i^0$ to $p_i^1$ is
$-\int_{p_i^0}^{p_i^1} x_i(t, m)\,\dd t$.
\end{definition}

The loss of \omterm{def:b1:consumer-welfare:consumer-surplus}{consumer surplus} from the doubled price in the figure is
$\int_1^2 500/t\,\dd t = 500\ln 2 = 346.57$, between EV and CV. That holds in
general for a \omterm{def:b1:consumer-choice:normal-inferior}{normal good}: the \omterm{def:b1:consumer-choice:demand-function}{Marshallian demand} through $(x_i(p^0), p^0)$
and $(x_i(p^1), p^1)$ lies between the two compensated curves, one through each
of the two points.

\begin{proposition}[Exact surplus with quasi-linear utility]\label{prop:b1:consumer-welfare:quasi}
\index{quasi-linear utility!consumer surplus}
If $u = \phi(x_1) + x_2$ with $p_2 = 1$ and an interior optimum, then
$\mathrm{CV} = \mathrm{EV} = $ the loss of \omterm{def:b1:consumer-welfare:consumer-surplus}{consumer surplus}.
\end{proposition}

\begin{proof}
At an interior optimum the demand for good 1 solves $\phi'(x_1) = p_1$, so it
does not depend on income: $x_1(p_1, m) = h_1(p_1, u)$ for every $u$. The three
integrals of \cref{prop:b1:consumer-welfare:areas} and of the definition are
then integrals of the same function between the same prices.
\end{proof}

How far apart can the three be in general? The gap comes from the \omterm{def:b1:income-and-substitution-effects:effects}{income
effect}, so it is small when the good takes a small share of the budget or
has a small income elasticity. A result known as Willig's bound (admitted
here) makes this precise: for a price change whose surplus loss is a
fraction $A$ of income, CV and EV differ from it by about
$\eta A/2$ in relative terms, with $\eta$ the income elasticity. A surplus
loss of $5\%$ of income with $\eta = 1$ puts CV and EV within about $2.5\%$
of it. That is why applied welfare economics, from \cref{ch:b1:partial-equilibrium-welfare}
on, uses areas under estimated demand curves.

\section{The true cost-of-living index}\label{sec:b1:consumer-welfare:col}

\begin{definition}[True cost-of-living index]\label{def:b1:consumer-welfare:col-index}
\index{true cost-of-living index}\index{cost-of-living index}
The \emph{true cost-of-living index} between prices $p^0$ and $p^1$, at the
reference utility $u$, is $I(p^0, p^1; u) = e(p^1, u)/e(p^0, u)$: the factor by
which income must rise to keep the consumer at utility $u$.
\end{definition}

A statistical office does not know $e$. It knows prices and the basket
$x^0$ bought in the base period, and computes the Laspeyres index
$L = p^1\cdot x^0/p^0\cdot x^0$; when it knows the current basket $x^1$, the
Paasche index $P = p^1\cdot x^1/p^0\cdot x^1$.

\begin{proposition}[The Laspeyres and Paasche bounds]\label{prop:b1:consumer-welfare:bounds}
\index{Laspeyres index!bound}\index{Paasche index!bound}
For a consumer who chooses $x^0$ at $(p^0, m^0)$ with utility $u^0$, and $x^1$ at
$(p^1, m^1)$ with utility $u^1$,
\[
I(p^0, p^1; u^0) \le L ,\qquad I(p^0, p^1; u^1) \ge P .
\]
\end{proposition}

\begin{proof}
$x^0$ reaches $u^0$, so the cheapest way to reach $u^0$ at $p^1$ costs at most
$p^1\cdot x^0$: $e(p^1, u^0) \le p^1\cdot x^0$. And $e(p^0, u^0) = p^0\cdot x^0$,
since $x^0$ was chosen at $p^0$. Dividing gives the first inequality. In the
same way $e(p^0, u^1) \le p^0\cdot x^1$ and $e(p^1, u^1) = p^1\cdot x^1$, which
gives the second.
\end{proof}

\begin{omfigure}
\begin{tikzpicture}
\begin{axis}[omrecord, grid=none, xlabel={$x_1$}, ylabel={$x_2$},
  xmin=0, xmax=90, ymin=0, ymax=160, width=0.55\linewidth, height=0.42\linewidth, samples=80]
  \addplot[demand, domain=16:90] {2500/x};
  \addplot[supply] coordinates {(0,150) (75,0)};
  \addplot[shifted] coordinates {(0,141.42) (70.71,0)};
  \addplot[guide] coordinates {(0,100) (100,0)};
  \addplot[only marks, mark=*, black] coordinates {(50,50) (35.36,70.71)};
  \node[font=\footnotesize, anchor=south west] at (axis cs:50,50) {$x^0$};
  \node[font=\footnotesize, anchor=south west] at (axis cs:35.36,70.71) {cheapest};
\end{axis}
\end{tikzpicture}

{\small $u = x_1x_2$, base prices $(1, 1)$ and income 100 (dotted line), base bundle
$x^0 = (50, 50)$. At new prices $(2, 1)$ the old bundle costs 150 (red line):
$L = 1.5$. The cheapest bundle on the old \omterm{def:b1:preferences-and-utility:indifference-curve}{indifference curve} costs 141.42
(dashed line): the true index is $1.4142$.}
\end{omfigure}

The Laspeyres index overstates the rise in the cost of living because the
consumer, compensated, would not keep buying $x^0$; it is the
substitution bias of the recall box, now proved. A chained index updates
the basket every period, so its error in each link is second order in that
period's price changes, and it stays closer to the true index. For
Cobb--Douglas preferences with budget shares $s_i$, the \omterm{def:b1:income-and-substitution-effects:expenditure}{expenditure
function} is $e(p, u) = u\,\prod_i p_i^{s_i}/\prod_i s_i^{s_i}$, so the true index
is the same at every reference utility and equals the geometric mean of
the price relatives, $I = \prod_i (p_i^1/p_i^0)^{s_i}$
(\cref{exo:b1:consumer-welfare:cd-index}). Since a weighted geometric mean is
below the weighted arithmetic mean, $L = \sum_i s_i\,p_i^1/p_i^0$ is above it,
as the proposition requires.

\begin{measure}[Fixed-basket and chained indices]
The CPI-U prices a basket of goods and services whose weights come from the
household expenditure surveys of \cref{ch:b1:consumer-choice}; since the
January 2023 data the weights are updated every year. Since July 2002 the
Bureau of Labor Statistics also publishes a chained index, the C-CPI-U,
with a superlative formula that lets the weights follow substitution across
categories of items. From 2001 to 2023 the CPI-U rose on average about
$0.2$ percentage points a year faster. % ledger: b1:ccpiu-gap
The deflator of the opening chart is chained too, but it also differs in
coverage and weights, so the gap of the chart, $2.61 - 2.16 = 0.45$ points a
year, is more than substitution bias. % figdata: us_growth
\end{measure}

\begin{model}{The cost-of-living index}
\assumes A representative household with stable preferences, whose
\omterm{def:b1:income-and-substitution-effects:expenditure}{expenditure function} defines the cost of reaching a fixed utility.
\predicts A fixed-basket index rises faster than the true index; a chained
index lies between them; the gap grows with the dispersion of price
changes across goods.
\tested The chained CPI has risen about 0.2 points a year more slowly than
the fixed-basket CPI-U since 2001. % ledger: b1:ccpiu-gap
A commission of five economists put the total upward bias of the American
CPI in the 1990s at 1.1 points a year, adding substitution, new goods
and quality change (Boskin and others 1998). % ledger: b1:boskin1998
\breaks Households differ: a household that cannot substitute, or that buys
more of the goods whose prices rose fastest, faces a different true index,
so no single index is true for all. New goods and quality changes need
values that no price survey records.
\end{model}

\begin{record}[A review of the CPI]
A commission of five economists estimated, and in 1998 reiterated, that the
American CPI overstated the rise in the cost of living by 1.1 percentage
points a year (Boskin and others 1998). % ledger: b1:boskin1998
The sources were substitution between goods (which a fixed basket misses),
substitution between outlets (discount stores), new products introduced
late, and quality improvements priced as price rises. The chained index of
the measure box addresses the first of these; the theory of this chapter
says it cannot address the others.
\end{record}

\section{New goods and quality}\label{sec:b1:consumer-welfare:new-goods}

A good that did not exist in the base period has no base price. The model
gives it one.

\begin{definition}[Virtual price]\label{def:b1:consumer-welfare:virtual-price}
\index{virtual price}\index{reservation price}
The \emph{virtual price} of a good that is not available is the price at
which the consumer's demand for it would be exactly zero, given her income
and the other prices: the intercept of its demand curve on the price axis.
\end{definition}

Pricing the good at its \omterm{def:b1:consumer-welfare:virtual-price}{virtual price} before it appears, the true index can
treat its arrival as a price fall from the \omterm{def:b1:consumer-welfare:virtual-price}{virtual price} to the market
price. The welfare gain is the area under the \omterm{def:b1:income-and-substitution-effects:compensated}{compensated demand} between
the two prices, or the consumer-surplus triangle when \omterm{def:b1:income-and-substitution-effects:effects}{income effects} are
small. A Laspeyres index that adds the good only years after it appears
misses this gain, and the falls in its price in the meantime.

\begin{omfigure}
\begin{tikzpicture}
\begin{axis}[omrecord, grid=none, xlabel={Quantity (gigabytes a year)}, ylabel={Price (dalens)},
  xmin=0, xmax=65, ymin=0, ymax=34, width=0.5\linewidth, height=0.4\linewidth]
  \addplot[cs area] coordinates {(0,30) (40,10) (0,10)} \closedcycle;
  \addplot[demand] coordinates {(0,30) (60,0)};
  \addplot[guide] coordinates {(0,10) (40,10) (40,0)};
  \addplot[only marks, mark=*, black] coordinates {(0,30) (40,10)};
  \node[font=\footnotesize, anchor=west] at (axis cs:1,31.5) {virtual price 30};
  \node[font=\footnotesize] at (axis cs:12,15) {400};
\end{axis}
\end{tikzpicture}

{\small Demand for phone data in Brenland (invented), $q = 60 - 2p$. Before the
service exists its \omterm{def:b1:consumer-welfare:virtual-price}{virtual price} is 30, where demand is zero. At a market
price of 10 a household buys 40 gigabytes, and the triangle, 400 dalens a
year, is the value of the new good.}
\end{omfigure}

\begin{record}[A new good]
Cellular telephones were introduced in the United States in 1983, and about
55 million were in use at the end of 1997, but the CPI did not include them
until 1998. % ledger: b1:hausman1999
Hausman (1999) estimated the demand for cellular service, valued its
arrival and its falling prices, and found that a correctly computed index
of telecommunications services would have fallen by about $0.8\%$ a year,
and that leaving the telephones out had biased that index upward by between
0.8 and 1.9 per cent a year. % ledger: b1:hausman1999
\end{record}

Quality changes raise the same problem in a milder form: a computer sold
at the same price as last year's but twice as fast is cheaper per unit of
computing. Statistical offices adjust some prices for quality, with
regressions of prices on characteristics, but most goods are adjusted
little or not at all; the commission's 1.1 points include its estimate of
what this misses.

\section*{Exercises}\label{sec:b1:consumer-welfare:exercises}

\begin{exercise}[$\star$]\label{exo:b1:consumer-welfare:cs-linear}
Demand is $q = 100 - 2p$. The price rises from 10 to 15. Compute the change
in \omterm{def:b1:consumer-welfare:consumer-surplus}{consumer surplus}.
\end{exercise}

\begin{exercise}[$\star$]\label{exo:b1:consumer-welfare:growth}
In the opening chart's table the CPI-U is 240.01 and the deflator 206.89 in
2024, both with 1990 = 100. % figdata: us_growth
Compute the average yearly growth of each over the 34 years.
\end{exercise}

\begin{exercise}[$\star$]\label{exo:b1:consumer-welfare:lp}
Two goods have prices $(1, 1)$ and quantities $(10, 10)$ in the base year, and
prices $(2, 1)$ and quantities $(6, 12)$ in the current year. Compute the
Laspeyres and Paasche indices.
\end{exercise}

\begin{exercise}[$\star$]\label{exo:b1:consumer-welfare:signs}
What are the signs of CV and EV for a fall in the price of a good the
consumer buys? Interpret each in words.
\end{exercise}

\begin{exercise}[$\star\star$]\label{exo:b1:consumer-welfare:cv-ev}
For $u = \sqrt{x_1x_2}$, income 1000 and prices $(1, 1)$, the price of good 1
doubles. Compute CV and EV from the \omterm{def:b1:income-and-substitution-effects:expenditure}{expenditure function}.
\end{exercise}

\begin{exercise}[$\star\star$]\label{exo:b1:consumer-welfare:between}
On the numbers of the previous exercise, compute the loss of \omterm{def:b1:consumer-welfare:consumer-surplus}{consumer
surplus} and check that it lies between EV and CV.
\end{exercise}

\begin{exercise}[$\star\star$]\label{exo:b1:consumer-welfare:cd-index}
Show that for Cobb--Douglas preferences the \omterm{def:b1:consumer-welfare:col-index}{true cost-of-living index} is
$\prod_i(p^1_i/p^0_i)^{s_i}$ at every reference utility. For the consumer of
the figure of the Laspeyres bound, compute it and the error of the
Laspeyres index.
\end{exercise}

\begin{exercise}[$\star\star$]\label{exo:b1:consumer-welfare:basket}
Brenland's 2015 basket costs 10,850 dalens, of which bread 600, clothing 800,
energy 750, rent 8,400 and transport 300; the price relatives 2015--2025 are
1.355, 1.225, 1.84, 1.4288 and 1.3033 (invented data). Rebuild the 2025
Laspeyres index and the geometric index, and compare with the chained index
of the weekend problem's table, 143.27. % figdata: brenland_indices
\end{exercise}

\begin{exercise}[$\star\star$]\label{exo:b1:consumer-welfare:spread}
For small price changes,
$\ln L - \ln I \approx \tfrac12\sum_i s_i(r_i - \bar r)^2$, where $r_i$ is the log
price relative of good $i$ and $\bar r = \sum_i s_ir_i$ (Cobb--Douglas true index).
Two goods with equal shares see their prices grow 1 point a year faster and
1 point a year slower than average. Predict the substitution bias of a
basket fixed for one year, and of one fixed for ten years.
\end{exercise}

\begin{exercise}[$\star\star\star$]\label{exo:b1:consumer-welfare:prove-bound}
Prove the two bounds of \cref{prop:b1:consumer-welfare:bounds}, and show by
an example with \omterm{def:b1:preferences-and-utility:perfect}{perfect complements} that the Laspeyres index can equal the
true index.
\end{exercise}

\begin{exercise}[$\star\star\star$]\label{exo:b1:consumer-welfare:quasi}
For $u = 10\sqrt{x_1} + x_2$, $p_2 = 1$ and income 100, the price of good 1 rises
from 1 to 2. Compute the demand, CV, EV and the loss of \omterm{def:b1:consumer-welfare:consumer-surplus}{consumer surplus},
and check \cref{prop:b1:consumer-welfare:quasi}.
\end{exercise}

\begin{exercise}[$\star\star\star$]\label{exo:b1:consumer-welfare:low-income}
``The CPI overstates the rise in the cost of living.'' What evidence would
refute this for a low-income household? Name one way its true index can
rise faster than the CPI.
\end{exercise}

\section{Problem: Brenland's Statistician Revises the CPI}\label{sec:b1:consumer-welfare:problem}

\begin{problem}[{Weekend problem --- the substitution bias of Brenland's CPI}]\label{pb:b1:consumer-welfare:cpi}
Brenland's CPI (invented data) prices a basket of five goods bought by a
household in 2015. The table below gives the basket, its 2015 cost and the
prices of 2025. Households have Cobb--Douglas preferences with the 2015
budget shares $s_i$.

\begin{center}
\pgfplotstabletypeset[col sep=comma, string type,
  columns={good,qty,p2015,p2025,spend2015,share,relative},
  columns/good/.style={column name=Good},
  columns/qty/.style={column name=Quantity},
  columns/p2015/.style={column name=Price 2015},
  columns/p2025/.style={column name=Price 2025},
  columns/spend2015/.style={column name=Cost 2015},
  columns/share/.style={column name=Share},
  columns/relative/.style={column name=Relative},
  every head row/.style={before row=\toprule, after row=\midrule},
  every last row/.style={after row=\bottomrule}]
  {figdata/bachelor-1/07-consumer-welfare/out/brenland_basket.csv}
\end{center}

\noindent The script's chained index for 2025 (2015 = 100) is 143.27; questions
2, 5 and 6 compute the other three. % figdata: brenland_indices

\medskip\noindent\textbf{Part I --- Two fixed baskets.}
\begin{enumerate}
  \item Check the 2015 cost, 10,850 dalens, and the shares of rent and energy.
  \item Compute the cost of the 2015 basket at 2025 prices and the Laspeyres index.
  \item Show that the Laspeyres index is the share-weighted arithmetic mean of the price relatives.
  \item Show that, with Cobb--Douglas households, the 2025 quantities are $s_i m^1/p^1_i$, and that the Paasche index is then $1/\sum_i s_i/(p^1_i/p^0_i)$.
  \item Compute the Paasche index.
  \item Compute the geometric index $\prod_i (p^1_i/p^0_i)^{s_i}$.
  \item Order the three, and say which mean each is.
\end{enumerate}

\medskip\noindent\textbf{Part II --- The true index and the bias.}
\begin{enumerate}[resume]
  \item Explain why the geometric index is the \omterm{def:b1:consumer-welfare:col-index}{true cost-of-living index} for these households, whatever the reference utility.
  \item Compute the substitution bias of the Laspeyres index in 2025, in index points.
  \item Compute the average yearly growth of the Laspeyres and the geometric index, and the bias per year.
  \item Explain why the bias is so small in Brenland.
  \item In 2022 the price of energy rose by $36.7\%$, the other four prices by
    $6.7\%$ to $8.3\%$. % computed: code/shared/country basket table, prices 2021 and 2022
    The script's table shows the Laspeyres--geometric gap at 0.04 points in
    2021 and 0.37 points in 2022. Why does a year like 2022 widen the
    gap? % figdata: brenland_indices
  \item Where does the chained index lie? Why does it not reach the geometric index?
\end{enumerate}

\medskip\noindent\textbf{Part III --- A new good.} Phone data arrived in
Brenland in 2020. Each household's demand is $q = 60 - 2p$ (gigabytes a year,
$p$ in dalens per gigabyte); the price was 20 in 2020 and 10 in 2025. Treat
utility as quasi-linear in other spending.
\begin{enumerate}[resume]
  \item What is the \omterm{def:b1:consumer-welfare:virtual-price}{virtual price} of phone data before 2020?
  \item Compute the quantity bought and the \omterm{def:b1:consumer-welfare:consumer-surplus}{consumer surplus} in 2020 and in 2025.
  \item In 2025, how much does the new good lower the cost of reaching the 2015 utility, in dalens and as a share of the 2025 cost of the 2015 basket?
  \item Compute the true index in 2025 with the new good, starting from the geometric index.
  \item Compute its average yearly growth over 2015--2025, and the bias per year of an index that leaves the phone out.
  \item The CPI adds phone data only in 2025, at its 2025 price. What does it miss?
\end{enumerate}

\medskip\noindent\textbf{Part IV --- A pension.} A pension of 15,000 dalens a
year in 2015 is indexed to the CPI. The script's table gives, for 2016 to
2025, Laspeyres--geometric gaps whose sum is 1.81 index points. % figdata: brenland_indices
\begin{enumerate}[resume]
  \item Compute the 2025 pension indexed to the Laspeyres index, and to the geometric index.
  \item How much more does the Laspeyres pension pay in 2025, and over 2016--2025 in all?
  \item Compute the 2025 pension indexed to the chained index. What would a switch to it cost the pensioner in 2025?
  \item Which index would keep a pensioner who uses no phone data exactly as well off as in 2015?
  \item A pensioner spends more of her budget on energy than the average household. Does the geometric index with average shares overstate or understate her true index? Explain.
  \item State the result: the substitution bias of Brenland's CPI over 2015--2025.
\end{enumerate}
\end{problem}
